\subsection{SM4}
\subsubsection{SM4简介}
SM4是中国国家密码管理局于2012年发布的分组密码算法，分组长度为128位，密钥长度为128位，采用32轮加密。
SM4的加解密步骤是一样的，只不过轮密钥顺序相反。

\subsubsection{轮函数}
先将明文分为4个32位分组，然后进行32轮加密，每轮使用一个32位的子密钥，生成一个新的32位分组。
$$
X_{i+4} = X_i \oplus T(X_{i+1} \oplus X_{i+2} \oplus X_{i+3} \oplus RK_i)
$$

合成置换T是一个可逆置换，由非线性变换S和线性变换L组成，S使用S盒进行以字节为单位的置换。

$$
S
\begin{bmatrix}
D6 & 90 & E9 & FE & CC & E1 & 3D & B7 & 16 & B6 & 14 & C2 & 28 & FB & 2C & 05 \\
2B & 67 & 9A & 76 & 2A & BE & 04 & C3 & AA & 44 & 13 & 26 & 49 & 86 & 06 & 99 \\
9C & 42 & 50 & F4 & 91 & EF & 98 & 7A & 33 & 54 & 0B & 43 & ED & CF & AC & 62 \\
E4 & B3 & 1C & A9 & C9 & 08 & E8 & 95 & 80 & DF & 94 & FA & 75 & 8F & 3F & A6 \\
47 & 07 & A7 & FC & F3 & 73 & 17 & BA & 83 & 59 & 3C & 19 & E6 & 85 & 4F & A8 \\
68 & 6B & 81 & B2 & 71 & 64 & DA & 8B & F8 & EB & 0F & 4B & 70 & 56 & 9D & 35 \\
1E & 24 & 0E & 5E & 63 & 58 & D1 & A2 & 25 & 22 & 7C & 3B & 01 & 21 & 78 & 87 \\
D4 & 00 & 46 & 57 & 9F & D3 & 27 & 52 & 4C & 36 & 02 & E7 & A0 & C4 & C8 & 9E \\
EA & BF & 8A & D2 & 40 & C7 & 38 & B5 & A3 & F7 & F2 & CE & F9 & 61 & 15 & A1 \\
E0 & AE & 5D & A4 & 9B & 34 & 1A & 55 & AD & 93 & 32 & 30 & F5 & 8C & B1 & E3 \\
1D & F6 & E2 & 2E & 82 & 66 & CA & 60 & C0 & 29 & 23 & AB & 0D & 53 & 4E & 6F \\
D5 & DB & 37 & 45 & DE & FD & 8E & 2F & 03 & FF & 6A & 72 & 6D & 6C & 5B & 51 \\
8D & 1B & AF & 92 & BB & DD & BC & 7F & 11 & D9 & 5C & 41 & 1F & 10 & 5A & D8 \\
0A & C1 & 31 & 88 & A5 & CD & 7B & BD & 2D & 74 & D0 & 12 & B8 & E5 & B4 & B0 \\
89 & 69 & 97 & 4A & 0C & 96 & 77 & 7E & 65 & B9 & F1 & 09 & C5 & 6E & C6 & 84 \\
18 & F0 & 7D & EC & 3A & DC & 4D & 20 & 79 & EE & 5F & 3E & D7 & CB & 39 & 48 \\
\end{bmatrix}
$$
$$
L(B) = B \oplus (B <<< 2) \oplus (B <<< 10) \oplus (B <<< 18) \oplus (B <<< 24)
$$

最终密文是倒序的4个32个分组拼接而成。

\subsubsection{密钥扩展}
SM4的密钥扩展先把密钥分为4个32位子密钥MK，再使用其与FK生成中间变量K，
$K_i = MK_i \oplus FK_i$。
$$
FK =
\begin{bmatrix}
\texttt{A3B1BAC6} & \texttt{56AA3350} & \texttt{677D9197} & \texttt{B27022DC}
\end{bmatrix}
$$
之后的步骤类似轮函数的操作，最后生成的轮密钥$RK_i=K_{i+4}$。

$$
K_{i+4} = K_i \oplus T'\left(K_{i+1} \oplus K_{i+2} \oplus K_{i+3} \oplus CK_i\right)
$$
$$
CK =
\begin{bmatrix}
\texttt{00070E15} & \texttt{1C232A31} & \texttt{383F464D} & \texttt{545B6269} \\
\texttt{70777E85} & \texttt{8C939AA1} & \texttt{A8AFB6BD} & \texttt{C4CBD2D9} \\
\texttt{E0E7EEF5} & \texttt{FC030A11} & \texttt{181F262D} & \texttt{343B4249} \\
\texttt{50575E65} & \texttt{6C737A81} & \texttt{888F969D} & \texttt{A4ABB2B9} \\
\texttt{C0C7CED5} & \texttt{DCE3EAF1} & \texttt{F8FF060D} & \texttt{141B2229} \\
\texttt{30373E45} & \texttt{4C535A61} & \texttt{686F767D} & \texttt{848B9299} \\
\texttt{A0A7AEB5} & \texttt{BCC3CAD1} & \texttt{D8DFE6ED} & \texttt{F4FB0209} \\
\texttt{10171E25} & \texttt{2C333A41} & \texttt{484F565D} & \texttt{646B7279}
\end{bmatrix}
$$
$$
T'(A) = L'(S(A)),\quad
L'(B) = B \oplus (B <<< 13) \oplus (B <<< 23)
$$

